%% Applied Data Analysis -- Syllabus, Fall 2026 (English)
%% Compile with: pdflatex syllabus_ADA_2026_en.tex
\documentclass[11pt]{article}

\usepackage[a4paper,top=2.2cm,bottom=2.2cm,left=2.2cm,right=2.2cm]{geometry}
\usepackage[T1]{fontenc}
\usepackage[utf8]{inputenc}
\usepackage{lmodern}
\usepackage{microtype}
\usepackage{array}
\usepackage{longtable}
\usepackage{booktabs}
\usepackage{enumitem}
\usepackage{parskip}
\usepackage{titlesec}
\usepackage{xcolor}
\usepackage[hidelinks]{hyperref}

\definecolor{oxfordblue}{HTML}{002147}

\titleformat{\section}{\large\bfseries\color{oxfordblue}}{}{0pt}{\MakeUppercase}
\titlespacing*{\section}{0pt}{12pt}{4pt}

\newcommand{\hdivider}{\noindent\rule{\textwidth}{0.8pt}}
\newcommand{\readings}[1]{\par\vspace{3pt}{\footnotesize\itshape #1\par}}

\setlist[itemize]{leftmargin=1.4em,topsep=2pt,itemsep=1pt}
\setlength{\LTpre}{4pt}
\setlength{\LTpost}{4pt}

\newcolumntype{W}{>{\raggedright\arraybackslash}p{0.8cm}}
\newcolumntype{T}{>{\raggedright\arraybackslash}p{10.6cm}}
\newcolumntype{D}{>{\raggedright\arraybackslash}p{4.1cm}}

\pagestyle{plain}

\begin{document}

%% ---------------------------------------------------------------- HEADER
\begin{center}
{\LARGE\bfseries\color{oxfordblue} APPLIED DATA ANALYSIS}\\[2pt]
{\large Faculty of National Economy}\\[1pt]
{\large Bratislava University of Economics and Business}\\[1pt]
{\large Department of Economic Policy}
\end{center}

\hdivider
\smallskip

\noindent
\begin{tabular}{@{}>{\bfseries}p{3.1cm}p{12.9cm}@{}}
Instructors: & Tomáš Oleš, Martin Vaško, Ella Sargsyan \\
Email: & tomas.oles@euba.sk; martin.vasko@euba.sk;\newline esargsyan@cerge-ei-foundation.org \\
Office Hours: & Wednesday: 1--3:00 PM \\
Office Location: & Old Building, 4B40 \\
Credit Hours: & 6 \\
Term: & Fall 2026 \\
\end{tabular}

\smallskip
\hdivider

%% ---------------------------------------------------------------- DESCRIPTION
\section{Course Description}

The first goal of the course is to teach students knowledge and skills in modern methods of applied data
analysis and statistical learning, including the use of R software, to conduct empirical economic research
and design research approaches and methods for solving economic problems.

The second aim of this course is to acquaint students with basic econometric concepts and methods which
are necessary to read and understand current empirical research in economics, as well as plan and execute
independent research projects. The course will introduce the idea of regression analysis and cover some of
the most recent econometric techniques central to modern econometric practice, which aims to identify
causal relationships.

The course is delivered in two parts. \textbf{Part I -- Applied Data Analysis in R} (Weeks 1--7) covers the
data-analytic toolkit: the R language, data visualization, data wrangling, text analysis, and an
introduction to statistical learning through classification and clustering. \textbf{Part II -- Econometrics
in a Nutshell} (Weeks 8--13) is delivered as a self-contained module by Ella Sargsyan, Ph.D.\
(CERGE-EI Foundation), and covers regression analysis and modern causal-inference techniques.

%% ---------------------------------------------------------------- OUTCOMES
\section{Student Learning Outcomes}

Upon completing the course, students should acquire knowledge of modern methods of data research and
visualization, classification, clustering, linear regression, and general data analysis. Skills in working
with data, which they will be able to apply in their own empirical research. They will also gain advanced
skills in using modern software (R) in empirical economic research and will be able to write their own
functions.

Students will know basic econometric concepts, basic estimation methods, and methods for testing statistical
hypotheses. They will be able to apply standard methods of constructing econometric models, process
statistical information, obtain statistically sound conclusions, and give meaningful interpretation to the
results of the estimated econometric models. In addition, students will gain real data processing skills,
using econometric packages for building and estimating regression models in R.

%% ---------------------------------------------------------------- PREREQS
\section{Prerequisites}

Introduction to Statistics.

%% ---------------------------------------------------------------- TEXTBOOKS
\section{Textbooks and Course Material}

\begin{itemize}
  \item Wickham, H., Çetinkaya-Rundel, M., Grolemund, G. (2023). \textit{R for Data Science}. O'Reilly
        Media, Inc. Freely available at \url{https://r4ds.hadley.nz/}.
  \item Gareth, J., Daniela, W., Trevor, H., Robert, T. (2013). \textit{An Introduction to Statistical
        Learning: With Applications in R}. Springer. Freely available at \url{https://www.statlearning.com/}.
  \item Silge, J., Robinson, D. (2017). \textit{Text Mining with R: A Tidy Approach}. O'Reilly Media, Inc.
        Freely available at \url{https://www.tidytextmining.com/}.
  \item Wooldridge, Jeffrey M. \textit{Introductory Econometrics: A Modern Approach}. 5th ed., US, 2012.
  \item Angrist, J. D., Pischke, J.-S. (2009). \textit{Mostly Harmless Econometrics: An Empiricist's
        Companion}. Princeton University Press.
  \item Tutorial on basic R: \url{https://cran.r-project.org/doc/contrib/Paradis-rdebuts_en.pdf}.
\end{itemize}

All course material (relevant textbook chapters, reading papers and lecture slides) will be uploaded on the
Moodle course webpage.

%% ---------------------------------------------------------------- GRADING
\section{Grading and Evaluation}

The final grade is made up of three components.

\begin{itemize}
  \item \textbf{Part I --- quizzes and home assignments: 30\%.} Short quizzes and assignments given during
        and after the lectures of Part I, together with the two home assignments based on R.
  \item \textbf{Part II --- Econometrics in a Nutshell: 35\%.} Assessed within that module.
  \item \textbf{Final data project: 35\%.} A data project presented at the end of the semester. It draws on
        both Part I and Part II and is the concluding piece of work for the course.
\end{itemize}

\noindent\textbf{In-class questions and discussions} are not graded, but active participation may be taken
into account at the instructor's discretion.

\smallskip
\noindent\textbf{Final grade calculation:} there is no final written exam. Students do not need to reach a
given threshold in each individual component; the final grade is based on the cumulative score.

%% ---------------------------------------------------------------- DATES
\section{Important Dates}

\begin{itemize}
  \item \textbf{Part I (Applied Data Analysis in R):} September 22 -- November 4, 2026
  \item \textbf{Part II (Econometrics in a Nutshell):} October 26 -- December 4, 2026
  \item \textbf{Final project presentations:} December 7 -- 11, 2026
  \item \textbf{Make-up week:} December 14 -- 18, 2026
\end{itemize}

\noindent\textbf{Class times.}

\smallskip
\noindent\begin{tabular}{@{}>{\bfseries}p{2.4cm}p{6.4cm}p{6.4cm}@{}}
\toprule
 & \textbf{Part I} --- Applied Data Analysis in R & \textbf{Part II} --- Econometrics in a Nutshell \\
\midrule
Course code & NND21254/21 (English class) & CERGE-EI module \\
Day & Wednesday, from 23 Sep 2026 & Lecture: Mon + Thu, from 26 Oct 2026\newline Exercise: Wednesday \\
Lecture & 11:00--12:30, room 4B26 & 11:00--12:30 CET, \textbf{online} \\
Exercise & 13:30--15:00, room D203 & 13:30--15:00, room D203 (in person) \\
\bottomrule
\end{tabular}

\smallskip
\noindent This syllabus states the times for the \textbf{English class}. The Slovak class
(NND21210/21) meets on Tuesdays --- see the Slovak version of this syllabus.

\smallskip
\noindent In Part II the lectures are delivered \textbf{online} by Ella Sargsyan, while the
exercise session remains \textbf{in person} in room D203 at the usual Wednesday time.

\smallskip
\noindent\footnotesize The instructor reserves the right to change the content of the course material if he
perceives a need due to postponement of class caused by inclement weather, instructor illness, etc., or due
to the pace of the course.\normalsize

%% ---------------------------------------------------------------- INTEGRITY
\section{Academic Integrity}

Academic integrity is essential to the pursuit of learning and scholarship in a university, and to ensuring
that a degree from the Bratislava University of Economics and Business is a strong signal of each student's individual
academic achievement. As a result, the University treats cases of cheating and plagiarism very seriously.
The Bratislava University of Economics and Business Ethical Code (\url{https://euba.sk/univerzita/eticky-kodex})
outlines the behaviors that constitute academic dishonesty and the processes for addressing academic
offenses. Potential offenses include, but are not limited to:

\smallskip
\noindent\textit{In papers and assignments:} Using someone else's ideas or words without appropriate
acknowledgment; submitting your own work in more than one course without the permission of the instructor;
making up sources or facts; obtaining or providing unauthorized assistance on any assignment.

\noindent\textit{On tests and exams:} Using or possessing unauthorized aids; looking at someone else's
answers during an exam or test; misrepresenting your identity / pretending to be someone else.

\noindent\textit{In academic work:} Falsifying documents or grades; falsifying or altering any documentation
required by the University, including (but not limited to) medical forms.

\smallskip
All suspected cases of academic dishonesty will be investigated following procedures outlined in the Ethical
Code. If students have questions or concerns about what constitutes appropriate academic behavior or
appropriate research and citation methods, they are expected to seek out additional information on academic
integrity from their instructors or from other institutional resources.

\clearpage

%% ---------------------------------------------------------------- SCHEDULE
\section{Class Schedule}

\begin{longtable}{@{}WTD@{}}
\toprule
\textbf{Wk} & \textbf{Topic Description} & \textbf{Dates \& Instructor} \\
\midrule
\endfirsthead
\toprule
\textbf{Wk} & \textbf{Topic Description} & \textbf{Dates \& Instructor} \\
\midrule
\endhead
\bottomrule
\endfoot

\multicolumn{3}{@{}l}{\textbf{\color{oxfordblue} PART I --- APPLIED DATA ANALYSIS IN R}}\\
\midrule

1 & \textbf{R basics.} Basic R syntax, data types, vectors arithmetic, indexing, sorting, sorting using
\texttt{dplyr}, and plotting using basic packages. GitHub.
\readings{Reading: Gareth et al. (2013), Ch.~1. Wickham et al. (2023), Ch.~1, 2, 4--6.}
& Wed 23 Sep 2026\newline Lecture 11:00--12:30, 4B26\newline Exercise 13:30--15:00, D203\newline \textit{T. Oleš} \\
\midrule

2 & \textbf{Visualization Part 1.} Data visualization principles, creating custom plots with
\texttt{ggplot2}, and studying the advantages and pitfalls of widely-used plots.
\readings{Reading: Wickham et al. (2023), Ch.~3, 7.}
& Wed 30 Sep 2026\newline Lecture 11:00--12:30, 4B26\newline Exercise 13:30--15:00, D203\newline \textit{T. Oleš} \\
\midrule

3 & \textbf{Visualization Part 2.} Spatial visualization principles, creating custom plots with
\texttt{tmap}. Spatial data overview (i.e.\ points, lines and polygons) and ending with a brief raster
example.
\readings{Reading: Wickham et al. (2023), Ch.~3, 7.}
& Wed 7 Oct 2026\newline Lecture 11:00--12:30, 4B26\newline Exercise 13:30--15:00, D203\newline \textit{T. Oleš} \\
\midrule

4 & \textbf{Data Summarizing, Tidying and Joining.} Importing data from different file formats, web
scraping, tidy data with \texttt{tidyverse}, wrangling data with \texttt{dplyr}, handling date and time
formats.
\readings{Reading: Wickham et al. (2023), Ch.~9--16.}
& Wed 14 Oct 2026\newline Lecture 11:00--12:30, 4B26\newline Exercise 13:30--15:00, D203\newline \textit{T. Oleš} \\
\midrule

5 & \textbf{Text Analysis.} The tidy text format and tokenization. Stop words and word frequencies.
Sentiment analysis with tidy data (Bing, AFINN, NRC lexicons). Term frequency, inverse document frequency
and \texttt{tf-idf}.
\readings{Reading: Silge and Robinson (2017), Ch.~1--3.}
& Wed 21 Oct 2026\newline Lecture 11:00--12:30, 4B26\newline Exercise 13:30--15:00, D203\newline \textit{T. Oleš} \\
\midrule

6--7\textsuperscript{**} & \textbf{Classification and Clustering.}\textsuperscript{**} Introduction to
qualitative classification. Classifiers: logistic regression, naive Bayes, K-nearest neighbors, generalized
linear models. Overfitting, confusion matrix, specificity and sensitivity. Unsupervised learning and the
notion of similarity: K-means, hierarchical clustering and dendrograms, choosing the number of clusters,
principal component analysis and its use in cluster analysis.
\readings{Reading: Gareth et al. (2013), Ch.~4 and Ch.~12.}
& Wed 28 Oct and 4 Nov 2026\newline Lecture 11:00--12:30, 4B26\newline Exercise 13:30--15:00, D203\newline \textit{T. Oleš, M. Vaško} \\
\midrule

\multicolumn{3}{@{}l}{\textbf{\color{oxfordblue} PART II --- ECONOMETRICS IN A NUTSHELL}}\\
\midrule

8 & \textbf{The Nature of Econometrics and Economic Data.} What is Econometrics and why do we need it?
Steps in empirical economic analysis. The structure of economics data. Causality and the notion of
``ceteris paribus'' in econometric data. Definition of simple regression model.
\readings{Reading: Wooldridge (2012), Ch.~1; Ch.~2 (2.1).}
& \footnotesize\textbf{Lecture (online)}\normalsize\newline Mon 26 Oct 11:00--12:30\newline Thu 29 Oct 11:00--12:30\newline \footnotesize\textbf{Exercise (in person)}\normalsize\newline Wed 28 Oct 13:30--15:00, D203\newline \textit{E. Sargsyan} \\
\midrule

9 & \textbf{Regression Analysis with Cross-Sectional Data: Part 1.} Ordinary Least Squares. Properties of
Ordinary Least Squares. Assumptions of the simple regression model.
\readings{Reading: Wooldridge (2012), Ch.~2 (2.2--2.5).}
& \footnotesize\textbf{Lecture (online)}\normalsize\newline Mon 2 Nov 11:00--12:30\newline Thu 5 Nov 11:00--12:30\newline \footnotesize\textbf{Exercise (in person)}\normalsize\newline Wed 4 Nov 13:30--15:00, D203\newline \textit{E. Sargsyan} \\
\midrule

10 & \textbf{Regression Analysis with Cross-Sectional Data: Part 2.} Multiple regression analysis.
Hypothesis testing. Reporting regression results.
\readings{Reading: Wooldridge (2012), Ch.~3; Ch.~4 (4.2--4.6).}
& \footnotesize\textbf{Lecture (online)}\normalsize\newline Mon 9 Nov 11:00--12:30\newline Thu 12 Nov 11:00--12:30\newline \footnotesize\textbf{Exercise (in person)}\normalsize\newline Wed 11 Nov 13:30--15:00, D203\newline \textit{E. Sargsyan} \\
\midrule

11 & \textbf{Regression Analysis: Further Topics and Issues.} Units of measurement and functional forms.
Regression analysis with qualitative information. Heteroskedasticity and serial correlation. Endogeneity:
omitted variables, selection, and reverse causality.
\readings{Reading: Wooldridge (2012), Ch.~2 (2.4); Ch.~3 (3.3); Ch.~7 (7.1--7.4); Ch.~8.}
& \footnotesize\textbf{Lecture (online)}\normalsize\newline Mon 16 Nov 11:00--12:30\newline Thu 19 Nov 11:00--12:30\newline \footnotesize\textbf{Exercise (in person)}\normalsize\newline Wed 18 Nov 13:30--15:00, D203\newline \textit{E. Sargsyan} \\
\midrule

12 & \textbf{Solving Endogeneity.} Instrumental Variables. Regression Discontinuity Design.
\readings{Reading: Wooldridge (2012), Ch.~15 (15.1--15.3); Angrist and Pischke (2009), Ch.~6 (pp.~189--195).
\newline Papers to be discussed: Acemoglu, D., \& Johnson, S. (2005). ``Unbundling institutions.''
\textit{Journal of Political Economy}, 113(5), 949--995. Carpenter, C., \& Dobkin, C. (2009). ``The effect
of alcohol consumption on mortality: regression discontinuity evidence from the minimum drinking age.''
\textit{American Economic Journal: Applied Economics}, 1(1), 164--182.}
& \footnotesize\textbf{Lecture (online)}\normalsize\newline Mon 23 Nov 11:00--12:30\newline Thu 26 Nov 11:00--12:30\newline \footnotesize\textbf{Exercise (in person)}\normalsize\newline Wed 25 Nov 13:30--15:00, D203\newline \textit{E. Sargsyan} \\
\midrule

13 & \textbf{Panel Data Methods.} Introducing time dimension. Fixed effects model.
Differences-in-Differences.
\readings{Reading: Wooldridge (2012), Ch.~14 (14.1); Angrist and Pischke (2009), Ch.~5 (pp.~169--182).
\newline Paper to be discussed: Card, D., \& Krueger, A. B. (1994). ``Minimum Wages and Employment: A Case
Study of the Fast-Food Industry in New Jersey and Pennsylvania.'' \textit{The American Economic Review},
84(4), 772--793.}
& \footnotesize\textbf{Lecture (online)}\normalsize\newline Mon 30 Nov 11:00--12:30\newline Thu 3 Dec 11:00--12:30\newline \footnotesize\textbf{Exercise (in person)}\normalsize\newline Wed 2 Dec 13:30--15:00, D203\newline \textit{E. Sargsyan} \\

\end{longtable}

\smallskip
\noindent\footnotesize\textsuperscript{**} Covered based on the strength of the group. These two sessions
run in parallel with the opening weeks of Part II.\normalsize

\vfill
\begin{center}
\itshape ``We are what we repeatedly do. Excellence, then, is not an act, but a habit.'' --- Aristotle
\end{center}

\end{document}
